\loai{Tính từ thông}
\begin{vidu}%[2P5-1.2B][Loai1]
Một khung dây phẳng có tiết diện $S =12 cm^2$ được đặt trong một từ trường đều, cảm ứng từ $B=5.10^{-2}$ T. Vector cảm ứng từ  hợp với pháp tuyến của mặt phẳng khung dây một góc $60\degree$. Tính độ lớn của từ thông qua tiết diện S.
\choice
{\True $3.10^{-5}$ Wb}
{$5.10^{-5}$ Wb}
{$2.10^{-5}$ Wb}
{$4.10^{-5}$ Wb}
\loigiai{Từ thông qua một vòng dây:
$\Phi =BS\cos \alpha =\left({{5.10}^{-2}}\right)\left({{12.10}^{-4}}\right)\cos 60\degree=3.10^{-5}\ Wb.$
}
\end{vidu}
\loai{Tính cảm ứng từ}
\begin{vidu}%[2P5-1.2K][Loai2]
Một khung dây phẳng có diện tích 10 $cm^2$ đặt trong từ trường đều, mặt phẳng khung dây hợp với đường cảm ứng từ một góc $30\degree$. Độ lớn từ thông qua khung là $3.10^{-5}$  Wb. Cảm ứng từ có giá trị là
\choice
{\True $B = 6.10^{-2}$ T}
{$B = 4.10^{-2}$ T}
{$B = 3.10^{-2}$ T}
{$B = 5.10^{-2}$ T}
\loigiai{
\begin{itemize}
\item Mặt phẳng khung dây hợp với đường cảm ứng từ một góc $30\degree \Rightarrow$  góc hợp bởi cảm ứng từ và vector pháp tuyến của mặt phẳng khung dây là $90^\circ - 30^\circ=60^\circ$.
\item Từ thông qua khung: $\Phi  = NBS\cos \alpha \Rightarrow B = \dfrac{\Phi}{NS\cos \alpha} = \dfrac{3.10^{-5}}{10.10^{-4}.\cos 60^\circ}=6.10^{-2}$ T.
\end{itemize}
}
\end{vidu}
\loai{Tính góc}
\begin{vidu}%[2P5-1.2K][Loai3]
Độ lớn của từ thông qua khung dây phẳng, kín có tiết diện $S = 100\ cm^2$ được đặt trong từ trường đều, cảm ứng từ $B=0,02$ T có giá trị $\Phi =\sqrt{3}.10^{-4}$ Wb. Xác định góc hợp bởi các đường sức từ với mặt phẳng khung dây.
\choice
{$30\degree$}
{\True $60\degree$}
{$90\degree$}
{$45\degree$}
\loigiai{
\begin{itemize}
\item Từ thông qua một vòng dây:
\begin{align*}
\Phi =BS\cos \alpha \Rightarrow \cos \alpha =\dfrac{\Phi }{BS}=\dfrac{\sqrt{3}.10^{-4}}{0,02.\left(100.10^{-4}\right)}=\dfrac{{\sqrt{3}}}{2}\Rightarrow \alpha =30\degree.
\end{align*}
\item Góc hợp bởi các đường sức từ với mặt phẳng khung dây:
$\beta =90\degree-30\degree=60\degree.$
\end{itemize}
}
\end{vidu}
\loai{Khung dây hình vuông}
\begin{vidu}%[2P5-1.2K][Loai6]
Một khung dây hình vuông có cạnh a gồm hai vòng dây. Khung được đặt trong một từ trường đều có cảm ứng từ $B = 0,08$ T sao cho các đường sức từ vuông góc với mặt phẳng khung. Biết từ thông qua khung là $4. 10^{-4}$ Wb. Giá trị của a là
\choice
{\True 5 cm}
{6,3 cm}
{2,5 mm}
{5 mm}
\loigiai{
Ta có: $\Phi = NBa^2 \Rightarrow a = 5$ cm}
\end{vidu}
\loai{Khung dây tam giác}
\begin{vidu}%[2P5-1.2K][Loai8]
Một đoạn dây dẫn có chiều dài 1 m được uốn thành một khung kín hình một tam giác đều. Khung dây được đặt trong một từ trường đều có cảm ứng từ $B = 0,1$ T sao cho các đường sức từ vuông góc với mặt phẳng khung dây. Từ thông qua khung có độ lớn là
\choice
{$9,6. 10^{-3}$ Wb}
{$5,6. 10^{-3}$ Wb}
{\True $4,8. 10^{-3}$ Wb}
{0,011 Wb}
\loigiai{
Chiều dài mỗi cạnh: $a=\dfrac{1}{3}\ m \Rightarrow S=\dfrac{a^2\sqrt{3}}{4}=\dfrac{\sqrt{3}}{36}\ m^2 \Rightarrow \Phi = BS = 4,8.10^{-3}$ Wb}
\end{vidu}
\loai{Tổng hợp}
\begin{vidu}%[2P5-1.2G][Loai9]
immini[thm]{
Hình tròn biểu diễn miền trong đó có từ trường đều, có cảm ứng từ B. Khung dây hình vuông cạnh a ngoại tiếp đường tròn. Công thức nào sau đây biểu diễn chính xác từ thông qua khung?
\choice
{$\pi Ba^2$ Wb}
{\True $\dfrac{\pi Ba^2}{4}$ Wb}
{$\dfrac{\pi a^2}{2B}$ Wb}
{$Ba^2$ Wb}
}{\begin{tikzpicture}[scale=1,thick,>=stealth']
\coordinate (o) at (0,0);
\draw [very thick,fill=white](-1.5,-1.5)--(-1.5,1.5)--(1.5,1.5)--(1.5,-1.5)--cycle ;
\draw [very thick,dashed](o)circle(1.5);
\path (o)--++(90:1)coordinate (b)node[right=6pt]{$\vec{B}$};
\draw (b)circle (0.2);
\filldraw[black] (b) circle(1.2pt);
\filldraw[black] (o) circle(1.2pt);
\end{tikzpicture}
}
\loigiai{
\begin{itemize}
\item Vì $\vec{B}$ vuông góc với mặt phẳng khung dây nên $\alpha = 0\degree$.
\item Trong đường tròn mới có từ trường nên diện tích S chính là diện tích hình tròn có bán kính $\dfrac{a}{2}$.
\item Khi đó, từ thông qua hình vuông được xác định bằng
\begin{align*}
\Phi =B\text{S}\cos 0\degree=B\pi \dfrac{a^2}{4}\ Wb.
\end{align*}
\end{itemize}
}
\end{vidu}
\loai{Thay đổi góc $\alpha$}
\begin{vidu}%[2P5-1.3K][Loai10]
Một khung dây chữ nhật ABCD gồm 20 vòng, cạnh là 5 cm và 4 cm. Khung đặt trong từ trường đều $B = 3.10^{-3}$ T, đường sức vuông góc với mặt khung. Quay khung $60\degree$ quanh cạnh AB. Tính độ biến thiên từ thông qua khung.
\choice
{\True $60.10^{-6}$ Wb}
{$50.10^{-6}$ Wb}
{$70.10^{-5}$ Wb}
{$30.10^{-5}$ Wb}
\loigiai{immini[thm]{
\begin{itemize}
\item Độ biến thiên từ thông:
\begin{align*}
\Delta \Phi  = {\Phi _{sau}} - {\Phi _{\text{trước}}} = NB{\rm{S}}(\cos {\alpha _{sau}} - \cos {\alpha _{\text{trước}}})
\end{align*}
\item Ban đầu, đường sức vuông góc với mặt phẳng khung $ \Rightarrow \alpha_{\text{trước}}=0^{\circ} $.
\item Lúc sau, khung quay đi một góc  $60^{\circ} \Rightarrow \alpha_{\text{sau}}=60^{\circ} $.
\item Thay các dữ kiện $\alpha_{sau} = 60\degree$, $\alpha_{\text{trước}}=0\degree$, $S=0,05.0,04=20.10^{-4} \ m^2$. Ta được:
\begin{align*}
\Delta \Phi  = \left|NBS(\cos {60\degree} - \cos {0\degree})\right|
\end{align*}
\begin{align*}
= \left|{20.3.10^{ - 3}}{.20.10^{ - 4}}.(0,5 - 1)\right| =  {60.10^{ - 6}} \ Wb.
\end{align*}
\end{itemize}
}{\begin{tikzpicture}[scale=1,thick,>=stealth']
\coordinate (o) at (0,0);
\draw (o)node[below left]{$A$}--++(150:4)--++(60:2)--++(-30:4)--cycle;
\draw [dashed](o)--++(90:4)coordinate(d)node[above left]{$D$}--++(60:2)coordinate(c)node[above right]{$C$}--++(270:4)coordinate(b2)node[below right]{$B$}--cycle;
\path(o)--++(150:2)--++(60:1)coordinate (n);
\draw [very thick,->] (n)--++(60:1.4)node[above right=-3pt]{\normalsize{\bth{\vec{n}}}}coordinate(n1);
\draw [very thick,->] (n)--++(0:1)node[below]{\normalsize{\bth{\vec{B}}}}coordinate(b);
\tkzMarkAngles[size=0.4cm](b,n,n1)
\path (n) ++(20:0.8) node{\normalsize{\bth{\alpha_{\text{sau}}}}};
\path(o)--++(90:2)--++(60:1)coordinate (m);
\path(m)--++(-90:0.1)coordinate (m1);
\draw [very thick,->] (m)--++(0:1.4)node[above right=-3pt]{\normalsize{\bth{\vec{n}}}};
\draw [very thick,->] (m1)--++(0:1)node[below]{\normalsize{\bth{\vec{B}}}}coordinate(b);
\filldraw[black] (n) circle(1.2pt);
\filldraw[black] (m) circle(1.2pt);
\filldraw[black] (m1) circle(1.2pt);
\draw [arrow data={0.4}{stealth}](o)++(90:4) arc (90:150:4) ;
\path(o)--++(60:2)coordinate(o1);
\draw [arrow data={0.4}{stealth}](o1)++(90:4) arc (90:150:4) ;
\end{tikzpicture}}
}
\end{vidu}
\loai{Thay đổi cảm ứng từ B}
\begin{vidu}%[2P5-1.3K][Loai11]
Một khung dây có diện tích $S=100 \ cm^2$ được đặt trong từ trường đều có $B = 0,2$ T. Vector pháp tuyến của mặt phẳng khung dây hợp với vector cảm ứng từ một góc $60\degree$. Ta thay đổi độ lớn cảm ứng từ $B’ = 0,1$ T. Độ biến thiên từ thông gửi qua khung dây là
\choice
{$4.10^{-4}$ Wb}
{$7.10^{-4}$ Wb}
{$6.10^{-4}$ Wb}
{\True $5.10^{-4}$ Wb}
\loigiai{
\begin{itemize}
\item Độ biến thiên từ thông gửi qua khung dây: $\Delta \Phi  = \Delta BS\cos \alpha$ 
\item Với: $\begin{cases}
\Delta B = \left|B_2-B_1\right| = |0,1-0,2|=0,1\ T\\
S = 100 \ cm^2\\
\alpha  = {60\degree}
\end{cases}\Rightarrow
\Delta \Phi = 5.10^{-4}\ Wb$.
\end{itemize}
}
\end{vidu}
\loai{Thay đổi diện tích S}
\begin{vidu}%[2P5-1.3K][Loai12]
Một khung dây có diện tích $S=100\ cm^2$  được đặt trong từ trường đều có $B = 0,2$ T. Pháp tuyến của  mặt phẳng khung dây hợp với vector cảm ứng từ một góc $60\degree$. Ta bóp méo khung dây sao cho diện tích của nó chỉ còn $S=25\ cm^2$. Độ biến thiên từ thông gửi qua khung dây là
\choice
{$6,5.10^{-4}$ Wb}
{\True $7,5.10^{-4}$ Wb}
{$8,5.10^{-5}$ Wb}
{$9,5.10^{-4}$ Wb}
\loigiai {
\begin{itemize}
\item Độ biến thiên từ thông gửi qua khung dây: $\Delta \Phi  = B\Delta S\cos \alpha$ 
\item Với: $\left\{ \begin{array}{l}
B = 0,2\ T\\
\Delta S = S_2 - S_1 = 75\ cm^2\\
\alpha  = 60\degree\\
\end{array} \right.$. Thay số:  $
\Delta \Phi = 7,5.10^{-4}\ Wb.
$
\end{itemize}
}
\end{vidu}
\loai{Đồ thị}
\begin{vidu}%[2P5-1.3K][Loai13]
immini[thm]{
Một khung dây cứng phẳng diện tích $25\ cm^2$ gồm 10 vòng dây, đặt trong từ trường đều, mặt phẳng khung vuông góc với các đường cảm ứng từ. Cảm ứng từ biến thiên theo thời gian như đồ thị hình vẽ. Tính độ biến thiên của từ thông qua khung dây kể từ $t = 0$ s đến $t = 0,4$\ s.
\choice
{$\Delta \Phi = 4.10^{-5}$\ Wb}
{$\Delta \Phi = 5.10^{-5}$\ Wb}
{\True $\Delta \Phi = 6.10^{-5}$\ Wb}
{$\Delta \Phi = 7.10^{-5}$\ Wb}}{\begin{tikzpicture}
\coordinate (o) at (0,0);
\draw [->](0,-0.3)--++(90:3)node[right]{\normalsize{\bth{B\;(T)}}};
\draw [->](-0.3,0)--++(0:2.5)node[above]{\normalsize{\bth{t\;(s)}}};
\draw[very thick](0,2)--(1.5,0);
\filldraw[black] (0,0)circle(1pt)node[below left]{\normalsize{\bth{O}}};
\filldraw[black] (1.5,0)circle(1pt)node[below]{\normalsize{\bth{0,4}}};
\filldraw[black] (0,2)circle(1pt)node[left]{\normalsize{\bth{2,4.10^{-3}}}};
\end{tikzpicture}}
\loigiai{
\begin{itemize}
\item Mặt phẳng khung vuông góc với các đường cảm ứng từ nên $\alpha =0\degree$.
\item Độ biến thiên từ thông
\begin{align*}
\Delta \Phi =\left|\Phi_2-\Phi_1\right|=\left|NS\cos \alpha (B_2-B_1)\right|=\left|10.25.10^{-4}\cos 0\degree.(0-2,4.10^{-3})\right|=6.10^{-5}\ Wb.
\end{align*}
\end{itemize}
}
\end{vidu}